% !TEX program = tectonic
% DOC MAÎTRE — « La donnée House : tout comprendre, tout vérifier, en 4 pages »
% Page 1 comprendre · Page 2 propre & complet · Page 3 backtestable + verdict + à-faire · Page 4 vérifié + carte.
% Tables + feux 🟢🟡🔴, quasi zéro prose. Chiffres = notebook Portefeuilles_House_Complet (vérifiés).
\documentclass[10pt,a4paper]{article}
\usepackage[margin=1.05cm]{geometry}
\usepackage[T1]{fontenc}
\usepackage[utf8]{inputenc}
\usepackage{lmodern}
\usepackage[french]{babel}
\usepackage{booktabs}
\usepackage{array}
\usepackage{xcolor}
\usepackage{amssymb}
\usepackage{tikz}
\usetikzlibrary{positioning}
\usepackage[most]{tcolorbox}

\definecolor{bleu}{RGB}{20,77,144}
\definecolor{vert}{RGB}{20,130,70}
\definecolor{orange}{RGB}{200,120,0}
\definecolor{rouge}{RGB}{175,35,35}
\definecolor{grisf}{RGB}{95,95,95}
\definecolor{grisc}{RGB}{205,205,205}

\setlength{\parindent}{0pt}\setlength{\parskip}{0pt}\pagestyle{empty}
\renewcommand{\arraystretch}{1.15}

\newcommand{\feu}[1]{\textcolor{#1}{\Large$\bullet$}}
\newcommand{\fv}{\feu{vert}}\newcommand{\fo}{\feu{orange}}\newcommand{\fr}{\feu{rouge}}
\newcommand{\tuile}[3]{\begin{tcolorbox}[nobeforeafter,width=0.237\textwidth,colback=#1!8,colframe=#1!55,
  boxrule=0.6pt,arc=2pt,left=1mm,right=1mm,top=1.2mm,bottom=1.2mm,halign=center,valign=center,height=17mm]
  {\large\bfseries\color{#1}#2}\\[0.3mm]{\scriptsize #3}\end{tcolorbox}}
\newtcolorbox{bloc}[2]{enhanced,breakable=false,colback=white,colframe=#1!55,boxrule=0.6pt,arc=2pt,
  left=2.2mm,right=2.2mm,top=1.2mm,bottom=1.4mm,toptitle=0.5mm,bottomtitle=0.5mm,boxsep=0pt,
  colbacktitle=#1,coltitle=white,fonttitle=\footnotesize\bfseries,title={#2},before skip=1.5mm,after skip=0mm}

\begin{document}

% ══════════════════════════════════ PAGE 1 — COMPRENDRE ══════════════════════════════════
\begin{center}
{\large\bfseries La donnée House --- tout comprendre, tout vérifier}\\[0.4mm]
{\scriptsize \textbf{le doc à ouvrir en premier} $\cdot$ House 2013--2026 $\cdot$ chiffres du notebook
\texttt{Portefeuilles\_House\_Complet}, recoupés cellule par cellule}
\end{center}
\vspace{0.6mm}

\begin{tcolorbox}[enhanced,colback=grisc!12,colframe=grisc,boxrule=0.5pt,arc=2pt,left=3mm,right=3mm,top=1.4mm,bottom=1.4mm]
{\footnotesize\textbf{Ce qu'on a fait.} Extraire \emph{proprement} toutes les transactions boursières
\emph{exploitables} des élus de la Chambre (2013--2026), et les reconstruire en portefeuilles quotidiens,
pour pouvoir \textbf{backtester}. Le trajet, en 2 actes / 9 étapes :}\\[1.4mm]
\begin{center}
\begin{tikzpicture}[x=0.98cm, y=0.9cm, font=\scriptsize]
\foreach \k/\lab/\c [count=\i from 0] in {%
  1/il arrive/bleu, 2/il trade/bleu, 3/il déclare/bleu, 4/il part/bleu,
  5/on rassemble/vert, 6/on trie/vert, 7/on valorise/vert, 8/on reconstruit/vert, 9/on vérifie/vert}{%
  \ifnum\k<9 \draw[\c!45, line width=0.9pt] ({\i*2.1+0.28}, 0) -- ({\i*2.1+1.82}, 0);\fi
  \fill[\c] ({\i*2.1}, 0) circle (0.23);
  \node[text=white, font=\scriptsize\bfseries] at ({\i*2.1}, 0) {\k};
  \node[text=\c, font=\scriptsize, anchor=north] at ({\i*2.1}, -0.30) {\lab};}
\node[bleu, font=\scriptsize\bfseries] at (2.05, 0.55) {ACTE I --- ils déposent};
\node[vert, font=\scriptsize\bfseries] at (12.6, 0.55) {ACTE II --- ce qu'on en fait};
\end{tikzpicture}
\end{center}
\vspace{0.6mm}
{\footnotesize Un dépôt de patrimoine porte \textbf{deux sections}, qui ne servent pas à la même chose :}\\[1mm]
\begin{center}
\begin{tikzpicture}[x=1cm, y=1cm, font=\scriptsize,
  bo/.style 2 args={rectangle, rounded corners=1.8pt, draw=#1, fill=#1!8, semithick,
                    inner sep=2.8pt, align=center, text width=#2, font=\scriptsize}]
\node[bo={grisf}{1.7cm}]  (m) at (0, 0)    {un \textbf{membre}};
\node[bo={bleu}{2.5cm}]   (d) at (3.0, 0)  {dépose des \textbf{documents}};
\node[bo={vert}{2.2cm}]   (l) at (6.2, 0)  {faits de \textbf{sections}};
\node[bo={orange}{6.0cm}] (p) at (11.7, 0.62) {\textbf{Schedule A} : ce qu'il possède $\rightarrow$ \textbf{positions} (le départ + le juge)};
\node[bo={orange}{6.0cm}] (t) at (11.7, -0.62) {\textbf{Schedule B} : ce qu'il achète/vend $\rightarrow$ \textbf{transactions} (le flux)};
\draw[->, grisf, semithick] (m) -- (d);
\draw[->, grisf, semithick] (d) -- (l);
\draw[->, grisf, semithick] (l.east) -- ++(0.4, 0) |- (p.west);
\draw[->, grisf, semithick] (l.east) -- ++(0.4, 0) |- (t.west);
\end{tikzpicture}
\end{center}
\end{tcolorbox}
\vspace{2mm}

\begin{center}{\footnotesize\bfseries CE QU'ON A, EN QUATRE CHIFFRES}\end{center}
\vspace{0.5mm}
\noindent
\tuile{vert}{129\,538}{transactions valorisées, 3 sources}\hfill
\tuile{bleu}{408 / 900}{membres qui tradent / au registre}\hfill
\tuile{orange}{609\,886}{lignes de patrimoine déclaré}\hfill
\tuile{rouge}{12$\times$}{empilé 104\,130 vs stock 8\,887 M\$}
\vspace{2.5mm}

\begin{center}
{\footnotesize De \textbf{900} membres et \textbf{6} types de dépôts, on tire un flux de \textbf{129\,538}
transactions exploitables chez \textbf{408} membres --- reconstruites jour après jour en portefeuilles,
puis \textbf{jugées} par des documents qui n'ont pas servi à les construire.}
\end{center}

\newpage

% ══════════════════════════════════ PAGE 2 — PROPRE & COMPLET ══════════════════════════════════
\begin{center}{\large\bfseries 1. Est-ce PROPRE et COMPLET ?}\hfill
{\scriptsize \fv~solide \quad \fo~limite connue \quad \fr~manque structurel}\end{center}
\vspace{1mm}

\begin{bloc}{bleu}{CE QU'ON A EXTRAIT --- le flux de transactions, par source}
\footnotesize
\begin{tabular}{@{}c l r l l@{}}
 & \textbf{source} & \textbf{trades} & \textbf{montant} & \textbf{publié après l'opération}\\
\midrule
\fv & déclarations rapides \emph{P} (fil de l'eau) & 107\,976 & fourchette, milieu & médiane \textbf{27 j}\\
\fo & rapports annuels \emph{O/A/T} (électronique) & 14\,084 & fourchette, milieu & médiane \textbf{374 j}\\
\fo & Schedule B scanné (océrisé) & 7\,478 & \textbf{forfait 8\,000\,\$} & médiane \textbf{393 j}\\
\midrule
 & \textbf{flux unifié F} & \textbf{129\,538} & \multicolumn{2}{l}{chez \textbf{408} membres $\cdot$ traçable ligne à ligne}\\
\end{tabular}\\[1mm]
{\scriptsize\color{grisf} Chaque ligne porte sa \textbf{date d'opération ET sa date de divulgation}
(\texttt{flux\_house\_2dates.csv}).}
\end{bloc}

\begin{bloc}{vert}{A-T-ON TOUT PRIS ? --- taux d'exploitation par type de document (le « recall »)}
\footnotesize
\begin{tabular}{@{}c l r r r@{}}
 & \textbf{type de document} & \textbf{exploité} & \textbf{lu, 0 ligne} & \textbf{papier jamais océrisé}\\
\midrule
\fv & \emph{P} --- déclaration rapide (transactions) & \textbf{100 \%} & --- & ---\\
\fo & \emph{O} --- rapport annuel & 74,7 \% & 450 & 721\\
\fr & \emph{T} / \emph{H} --- photos entrée/sortie & $\sim$65 \% & \multicolumn{2}{l}{lues 0 ligne + papier}\\
\fr & \emph{C} --- photo de candidat & \textbf{39,7 \%} & 4\,515 & 1\,583\\
\end{tabular}\\[1mm]
{\scriptsize\color{grisf} \textbf{La transaction est prise à 100 \%.} Ce qui se perd, c'est la \emph{photo}
du patrimoine (point de départ) --- surtout du papier jamais océrisé.}
\end{bloc}

\begin{bloc}{grisf}{POURQUOI ÇA MANQUE --- réalité (jamais déposé) / source (illisible) / nous (perdu)}
\footnotesize
\begin{tabular}{@{}l r r r r@{}}
 \textbf{le manque} & \textbf{membres} & \textcolor{grisf}{\textbf{RÉALITÉ}} & \textcolor{orange}{\textbf{SOURCE}} & \textcolor{rouge}{\textbf{NOUS}}\\
 & & rien déposé & déposé, illisible & exploité puis perdu\\
\midrule
photo d'entrée absente & 142 & 9 & 102 & \textbf{31}\\
photo de sortie absente & 96 & 86 & 9 & 1\\
aucun rapport annuel couvert & 79 & 66 & 13 & 0\\
\end{tabular}\\[1mm]
{\scriptsize\color{grisf} L'immense majorité des manques = \textbf{jamais déposé} ou \textbf{papier illisible} :
on n'a laissé sur la table quasiment rien d'\emph{automatiquement} extractible.}
\end{bloc}

\newpage

% ══════════════════════════════════ PAGE 3 — BACKTESTABLE + VERDICT + À-FAIRE ══════════════════════════════════
\begin{center}{\large\bfseries 2. Est-ce BACKTESTABLE ?}\end{center}
\vspace{1mm}

\begin{bloc}{bleu}{FEU PAR DIMENSION}
\footnotesize
\begin{tabular}{@{}c l p{10.6cm}@{}}
\fv & \textbf{le flux de transactions} & complet, traçable, décomposable ligne à ligne, chez 408 membres.\\[0.6mm]
\fo & \textbf{le point de départ} & 166 des 408 membres partent d'un portefeuille connu ; les \textbf{242} autres démarrent de zéro.\\[0.6mm]
\fv & \textbf{l'horloge (dates)} & reconstruit à la \textbf{date d'opération} = borne haute. La \textbf{date de divulgation est portée} dans le flux $\Rightarrow$ le rejeu réaliste (27 j / 374 j plus tard) se lance en une ligne, quand tu veux.\\[0.6mm]
\fo & \textbf{les prix} & 12,1 \% des lignes tombent faute de cours (1\,234 tickers) --- \textbf{confirmé réel} (re-test 22/07).\\[0.6mm]
\fo & \textbf{l'échantillon} & 11 années civiles $\Rightarrow$ écart minimal détectable \textbf{4,6 \%/an}.\\
\end{tabular}
\end{bloc}

\begin{bloc}{vert}{VERDICT}
\footnotesize
\textbf{L'extraction est propre, vérifiée et essentiellement complète pour tout ce qui est automatiquement
extractible.} Le flux de \textbf{129\,538 transactions} (3 sources, 408 membres, 3 défauts corrigés, recoupé
au notebook, jugé à 9\,$\rightarrow$\,57~\% par un juge indépendant) \textbf{EST backtestable dès maintenant}.
Le test « réaliste » (rejeu à la publication) est \textbf{déjà outillé} --- les deux dates sont dans le flux ---
mais pas encore joué, \emph{par choix}. Les manques restants sont soit \textbf{structurels} (jamais déposé /
papier illisible), soit des \textbf{limites de méthode connues}, pas des erreurs.
\end{bloc}

\begin{bloc}{orange}{À FAIRE --- pour extraire plus / rendre le backtest plus réaliste}
\footnotesize
\begin{tabular}{@{}r l l@{}}
\textbf{1} & rejeu à la date de \textbf{divulgation} (test réaliste) & \emph{données prêtes} --- \texttt{flux\_house\_2dates.csv}\\
\textbf{2} & \textbf{Schedule A} des rapports annuels comme départs & 460 membres exploitables vs 243 aujourd'hui\\
\textbf{3} & récupérer les lignes \textbf{sans-ticker} récupérables & 10\,416 + 59\,948 sur 337\,593\\
\textbf{4} & autre \textbf{source de cours} pour les 1\,455 tickers morts & ou l'accepter comme limite\\
\textbf{5} & \emph{(optionnel, risqué)} gisement \textbf{OCR} des rapides & +2\,510 trades, mais 81 \% des symboles viennent d'un LLM\\
\end{tabular}
\end{bloc}

\newpage

% ══════════════════════════════════ PAGE 4 — VÉRIFIÉ + CARTE DES DOCS ══════════════════════════════════
\begin{center}{\large\bfseries 3. Pourquoi croire ces chiffres, et où trouver le reste}\end{center}
\vspace{1mm}

\begin{bloc}{vert}{COMMENT ON SAIT QUE C'EST JUSTE}
\footnotesize
\begin{tabular}{@{}c p{15.6cm}@{}}
\fv & les 5 tables de départ sont \textbf{gelées} (empreinte sha256 conforme au manifeste) --- rien n'a bougé sous nos pieds.\\
\fv & la population est \textbf{re-croisée} avec le référentiel officiel des mandats : \textbf{98,9 \%} de concordance.\\
\fv & \textbf{chaque cascade et chaque tableau croisé} sont vérifiés par un \texttt{assert} (la somme doit boucler, sinon le notebook s'arrête).\\
\fv & la \textbf{classe} d'un actif (action / fonds / monétaire) est décidée \textbf{avant} tout prix --- pas de fonds déguisé en action.\\
\fv & la reconstruction est \textbf{jugée} par des positions déclarées qui n'ont \emph{jamais} servi à la construire : rappel médian \textbf{9\,$\rightarrow$\,57~\%}, + 5 recomptes à la main (écarts $10^{-9}$ à $10^{-15}$).\\
\fv & \textbf{aucun chiffre n'est écrit en dur} : tout se régénère en ré-exécutant le notebook.\\
\end{tabular}
\end{bloc}

\begin{bloc}{bleu}{OÙ TROUVER QUOI --- les 4 documents de cette donnée}
\footnotesize
\begin{tabular}{@{}l l l@{}}
\textbf{document} & \textbf{pages} & \textbf{à quoi ça sert}\\
\midrule
\textbf{ce doc} (\texttt{SYNTHESE\_EXTRACTION}) & 4 & \textbf{l'entrée unique} : comprendre + être sûre que tout est bon\\
la fiche (\texttt{FICHE\_DONNEES\_HOUSE}) & 5 & \textbf{comprendre en détail}, étape par étape\\
le deck (\texttt{SLIDES\_DONNEE\_HOUSE\_court}) & 39 & \textbf{présenter} à un public\\
le notebook (\texttt{Portefeuilles\_House\_Complet}) & --- & \textbf{la preuve} : tous les chiffres, rejouable $\cdot$ produit \texttt{flux\_house\_2dates.csv}\\
\end{tabular}\\[1.5mm]
{\scriptsize\color{grisf} L'ancien deck long (80 p.) est \textbf{archivé} : le deck court 39 p. le remplace.}
\end{bloc}

\vfill
\begin{center}{\scriptsize\color{grisf} Reste à faire de ton côté : lancer la cellule \texttt{add2dates2026} du notebook
(elle écrit \texttt{flux\_house\_2dates.csv} et s'auto-vérifie).}\end{center}

\end{document}
